\documentclass[10pt, a4paper]{article}

\usepackage{homework}

\title{\textbf{《数值代数》上机作业12}}
\author{岳瀚阳\quad\quad\quad\quad 3230104443}
\date{\today}

\begin{document}

\maketitle

%----------------------------------------------------------------------
\section{幂法求多项式方程的模最大根}
%----------------------------------------------------------------------

对于首一多项式方程
\[
f(x) = x^n + a_{n-1}x^{n-1} + \cdots + a_1 x + a_0 = 0,
\]
其全部根 $x_1,\ldots,x_n$ 的模最大者即为所求。将求根问题转化为特征值问题，构造友矩阵（companion matrix）
\[
C = \begin{bmatrix}
0 & 0 & \cdots & 0 & -a_0 \\
1 & 0 & \cdots & 0 & -a_1 \\
0 & 1 & \cdots & 0 & -a_2 \\
\vdots & \vdots & \ddots & \vdots & \vdots \\
0 & 0 & \cdots & 1 & -a_{n-1}
\end{bmatrix}_{n\times n}.
\]
$C$ 的特征多项式为 $\det(\lambda I - C) = \lambda^n + a_{n-1}\lambda^{n-1} + \cdots + a_0$，故 $C$ 的特征值恰好为 $f(x)=0$ 的全部根。于是求模最大根等价于求 $C$ 的模最大特征值。

对 $C$ 应用幂法：任取非零初始向量 $v^{(0)}$，迭代
\[
v^{(k+1)} = \frac{C v^{(k)}}{\|C v^{(k)}\|_2}, \qquad
\lambda^{(k+1)} = \bigl(v^{(k+1)}\bigr)^{\!\top} C v^{(k+1)},
\]
其中 $\lambda^{(k)}$ 为 Rayleigh 商，用于近似模最大特征值。停止准则取 $|\lambda^{(k+1)} - \lambda^{(k)}| < 10^{-12}$，最大迭代次数 $5000$。

幂法的收敛速度取决于次大模特征值与最大模特征值之比
\[
\left|\frac{\lambda_2}{\lambda_1}\right|.
\]
该比值越小，收敛越快；若存在多个不同特征值其模均等于 $|\lambda_1|$（例如共轭复根），则标准幂法不收敛，需改用其他技巧。

%----------------------------------------------------------------------
\section{数值结果}
%----------------------------------------------------------------------

分别对以下三个方程运行幂法子程序，结果汇总于表~\ref{tab:results}。

\begin{table}[H]
    \centering
    \caption{幂法求三个方程的模最大根}
    \label{tab:results}
    \begin{tabular}{clcc}
        \toprule
        题号 & 方程 & 模最大根 & 迭代次数 \\
        \midrule
        (i)   & $x^3 + x^2 - 5x + 3 = 0$ & $-3.0000000000$ & 31 \\
        (ii)  & $x^3 - 3x - 1 = 0$       & $1.8793852416$  & 125 \\
        (iii) & 8次方程                   & $-100.0000000000$ & 15 \\
        \bottomrule
    \end{tabular}
\end{table}

以下逐题给出分析与验证。

%----------------------------------------------------------------------
\subsection{方程 (i): $x^3 + x^2 - 5x + 3 = 0$}

该多项式可因式分解为
\[
x^3 + x^2 - 5x + 3 = (x-1)^2(x+3),
\]
故精确根为 $x_1 = -3,\; x_2 = x_3 = 1$。模最大根为 $-3$，且唯一（重根 $1$ 的模远小于 $3$）。幂法结果 $-3.0000000000$，迭代 31 次收敛，验证 $|f(-3)| \approx 2.87\times10^{-12}$。

收敛较快的原因在于次大模特征值与最大模特征值之比为 $|1|/|-3| = 1/3 \approx 0.333$，几何收敛因子小。

%----------------------------------------------------------------------
\subsection{方程 (ii): $x^3 - 3x - 1 = 0$}

令 $x = 2\cos\theta$，代入得
\[
8\cos^3\theta - 6\cos\theta - 1 = 2\cos 3\theta - 1 = 0
\;\Longrightarrow\; \cos 3\theta = \frac{1}{2}.
\]
故三个实根为
\[
x_1 = 2\cos\frac{\pi}{9} \approx 1.87938524, \quad
x_2 = 2\cos\frac{7\pi}{9} \approx -1.53208889, \quad
x_3 = 2\cos\frac{13\pi}{9} \approx -0.34729636.
\]
模最大根为 $x_1 = 2\cos(\pi/9)$，幂法给出 $1.8793852416$，误差约 $3.9\times10^{-13}$，验证 $|f(x_1)| \approx 2.97\times10^{-12}$。

本题迭代次数（125次）明显多于方程 (i)，原因是次大模根与最大模根之比为
\[
\frac{|x_2|}{|x_1|} = \frac{1.53208889}{1.87938524} \approx 0.815,
\]
接近 $1$，幂法收敛极慢。理论估计：达到 $10^{-12}$ 精度约需
\[
k \approx \frac{\ln 10^{-12}}{\ln 0.815} \approx 140
\]
次迭代，实测 125 次，基本吻合。

%----------------------------------------------------------------------
\subsection{方程 (iii): 8次方程}

方程为
\begin{multline*}
x^8 + 101x^7 + 208.01x^6 + 10891.01x^5 + 9802.08x^4 \\
+ 79108.9x^3 - 99902x^2 + 790x - 1000 = 0.
\end{multline*}

用 \texttt{numpy.roots} 求得全部根为
\[
-100,\; \pm 10\mathrm{i},\; -1\pm 3\mathrm{i},\; 1,\; \pm 0.1\mathrm{i}.
\]
各根的模分别为 $100,\; 10,\; 10,\; \sqrt{10}\approx 3.16,\; 1,\; 0.1$。模最大根为 $-100$，唯一且与次大模根之比为 $10/100 = 0.1$。幂法仅迭代 15 次即收敛到 $-100.0000000000$。

验证时直接代入计算 $|f(-100)| \approx 2.24$ 似乎偏大，但这源于浮点求值中的灾难性抵消：多项式在 $x=-100$ 处的首两项分别为 $(-100)^8 = 10^{16}$ 与 $101(-100)^7 = -1.01\times10^{16}$，数量级相近、符号相反，大量有效数字在相消中损失。若用 Horner 法或扩展精度重新估计，相对误差仍在机器精度量级。

%----------------------------------------------------------------------
\section{结果分析}
%----------------------------------------------------------------------

\textbf{（1）幂法的收敛速度严格取决于特征值模的分离程度。}
三题的分离比分别约为 $0.333$、$0.815$、$0.1$，对应迭代次数 $31$、$125$、$15$。分离比越小收敛越快；当分离比接近 $1$ 时（如方程 ii），即使算法正确，也需要上百次迭代才能达到双精度要求。

\textbf{（2）重根不影响幂法对模最大根的收敛。}
方程 (i) 中 $x=1$ 为二重根，但友矩阵对应特征值 $1$ 的几何重数为 $1$（Jordan 块结构），其模仍小于主导特征值 $-3$，幂法依旧收敛到正确根。

\textbf{（3）模最大根为实根时幂法给出实数结果，为复根时需采用复数运算。}
本题三题的模最大根均为实数（$-3$、$2\cos(\pi/9)$、$-100$），标准实幂法即可胜任。若模最大根为复数且存在其共轭，则需使用同时迭代或移位技巧。

%----------------------------------------------------------------------
\appendix
\section{源代码}
%----------------------------------------------------------------------

\subsection{\texttt{hw12\_test.py}（主程序）}
\lstinputlisting{../tests/hw12_test.py}

\subsection{\texttt{power\_method\_polynomial.py}（幂法子程序）}
\lstinputlisting{../algorithms/power_method_polynomial.py}

\end{document}
